\section{The forest estimate}\label{sec:forest}

\begin{theorem}[Forest estimate, KNT~\cite{KNT}]\label{thm:forest}
There is $r_0$ such that the following holds for all $r\ge r_0$. Let $G$ be a connected simple graph
of rank $r$. Let $W\subseteq G$ be a nonempty union of cycles with at most $r/\log^{2}r$ edges and
at most $(\log r)/100$ components. Then
\begin{equation}\label{eq:forest}
  p(G,W)\;\le\;\frac12+\frac{20}{\lgk{3}r}.
\end{equation}
\end{theorem}

This paper uses the theorem for a cycle $W$ only. The form with several components has the same
proof, and~\cite{Pan} needs it. The proof is that of~\cite{KNT}.

\begin{idea}{the proof in four steps}
\begin{enumerate}[leftmargin=1.5em]
\item We pass to a multigraph $\Gamma$ with all degrees at most $3$, in which $W$ has become a
small set $P$ of vertices (Lemma~\ref{lem:reduce}).
\item If $\Gamma-P$ has many disjoint short cycles, the estimate follows from
Lemma~\ref{lem:smallcuts}. So we may delete a small maximal family of them, and a graph $J$ of
large girth remains.
\item The graph $J$ has many cycles with between $6\log r$ and $\log r\sqrt{\log\log r}$ edges.
Their total weight $\sum_C2^{-|C|}$ is at least $(\frac18-o(1))\ln\log\log r$
(Lemma~\ref{lem:weight}).
\item For such a cycle $C$, the event that $C$ is a component of $\xi$ has probability at least
$2^{-|C|}$, and then the part of $\xi$ outside $P$ is not a forest. Most pairs of these events are
correlated by a factor of at most two (Lemma~\ref{lem:link}).
\end{enumerate}
With the factor two of step~4, the second moment of the number $Z$ of these events is at most
$2\mu^{2}+K\mu$, and Lemma~\ref{lem:second} then gives $\PP(Z=0)\le\half+K/(4\mu)$.
\end{idea}

We state the lemmas for the steps~3 and~4 with free parameters, because Section~\ref{sec:half}
uses them again.

\subsection{Reduction to a subcubic multigraph}

\begin{lemma}\label{lem:reduce}
Let $G$ and $W$ be as in Theorem~\ref{thm:forest}, with $m=|E(W)|$ and $c=\comp(W)$. There are a
connected multigraph $\Gamma$ with $\beta(\Gamma)=r$ and all degrees at most $3$, and a set
$P\subseteq V(\Gamma)$, such that
\begin{enumerate}
\item[\textup{(i)}] every vertex outside $P$ has degree $3$, $|P|\le2m$, and $\Gamma[P]$ has at
most $c$ components;
\item[\textup{(ii)}] $p(G,W)\le q(\Gamma,P)$.
\end{enumerate}
In particular $2r-2\le|V(\Gamma)|\le2r+2m$.
\end{lemma}

\begin{proof}
\step{Splitting.} Replace every vertex $v$ of degree $d\ge4$ by a path $v^{1}v^{2}\cdots v^{d-2}$.
Distribute the $d$ edges at $v$ along this path in some order: two to $v^{1}$, one to each of
$v^{2},\dots,v^{d-3}$, and two to $v^{d-2}$. If $v\in V(W)$, let the edges of $W$ at $v$ come first
in this order, so that they end in an initial segment $v^{1},\dots,v^{i}$ of the path with
$i<\deg_W(v)$. Let $G'$ be the new graph; its degrees are at most $3$. Let $P$ consist of these
initial segments and of the vertices of $W$ that were not split. Then $|P|\le\sum_v\deg_W(v)=2m$,
every edge of $W$ joins two vertices of $P$, and $G'[P]$ has at most $c$ components.

Contracting the new path edges maps $\Z(G')$ bijectively onto $\Z(G)$, so we may take
$\xi_G=\xi_{G'}\cap E(G)$. Every edge of $G'-P$ is a new path edge or an edge of $E(G)\sm E(W)$. If
$\xi_G\sm E(W)$ is a forest, then so is its union with the set of all new path edges, because
contracting these edges, which form a forest, gives back $\xi_G\sm E(W)$. Hence
$p(G,W)\le q(G',P)$.

\step{Cleaning.} As long as possible, delete a vertex of degree one outside $P$, or replace a vertex
of degree two outside $P$ and its two edges by a single edge. An even subgraph avoids a pendant edge
and contains both or none of the two edges at a vertex of degree two. So these operations induce
bijections of the cycle spaces. Let $\Gamma$ be the final multigraph. Each edge of $\Gamma-P$ is a
path of $G'$ all of whose vertices lie outside $P$. So if $\xi_\Gamma$ contains a cycle of
$\Gamma-P$, then the corresponding $\xi_{G'}$ contains a cycle of $G'-P$. This proves~(ii).

Edges with both ends in $P$ are never removed. So the vertices of $P$ keep degree at least two, and
the number of components of the graph induced on $P$ does not increase. The vertices outside $P$
now have degree exactly $3$. Finally $3|V(\Gamma)|-|P|\le2|E(\Gamma)|\le3|V(\Gamma)|$ and
$|E(\Gamma)|=|V(\Gamma)|+r-1$ give the bounds on $|V(\Gamma)|$.
\end{proof}

\subsection{Cycles of moderate length in a graph of large girth}\label{sec:moderate}

In this subsection $\Gamma$ is a connected multigraph with all degrees at most $3$, and $U$ is a
set of its vertices such that every vertex outside $U$ has degree $3$. We suppose that
$J=\Gamma-U$ is a simple graph of girth greater than $D$, where $D\ge3$ is an integer, and we put
$s=\lfloor(D-1)/2\rfloor$. Let $L\ge1$ be an integer, and let $\mathcal{C}$ be the set of the
cycles of $J$ of length at most $L$. Let $v$ be the number of vertices of $J$, and let
\[
  d_J=\sum_{u\in V(J)}\bigl(3-\deg_Ju\bigr)
\]
be its \emph{deficit}. This is the number of edges between $V(J)$ and $U$, so $d_J\le3|U|$.

A walk $u_0u_1\cdots u_\ell$ is \emph{non-backtracking} if $u_{i+1}\ne u_{i-1}$ for $0<i<\ell$. Let
$N_\ell(u,w)$ be the number of non-backtracking walks of length $\ell$ from $u$ to $w$ in $J$.

\begin{lemma}\label{lem:walks}
If $\ell>s$, then $N_\ell(u,w)\le3\cdot2^{\ell-s-1}$ for all $u,w$. Consequently, for every vertex
$w$ of $J$,
\begin{equation}\label{eq:load}
  \sum_{C\in\mathcal{C},\ C\ni w}2^{-|C|}\;\le\;L\,2^{-s}.
\end{equation}
\end{lemma}

\begin{proof}
There are at most $3\cdot2^{\ell-s-1}$ ways to choose the first $\ell-s$ steps. The last $s$ steps
form a non-backtracking walk of length $s$ between two given vertices. Such a walk is a path,
because $s$ is less than the girth. There is at most one such path, because two of them would give
a cycle of length at most $2s<D$. For the second statement, a cycle of length $\ell$ through $w$
can be traversed from $w$ in two directions. So there are at most $\half N_\ell(w,w)$ of them, and
$\ell>D>s$. Hence the cycles of each length $\ell\le L$ contribute at most $\frac34\,2^{-s}$.
\end{proof}

Let $B$ be the non-backtracking matrix of $J$. Its rows and columns are indexed by the directed
edges $(i\to j)$ of $J$, two for each edge. The entry in row $(i\to j)$ and column $(j'\to k)$ is
$1$ if $j'=j$ and $k\ne i$, and it is $0$ otherwise.

\begin{lemma}\label{lem:spectrum}
Suppose that $d_J<v$. Then, for every even $\ell\ge2$,
\[
  \tr B^{\ell}\;\ge\;\bigl(2-d_J/v\bigr)^{\ell}-3v\,2^{\ell/2}.
\]
\end{lemma}

\begin{proof}
Let $A$ be the adjacency matrix of $J$ and $\Delta$ the diagonal matrix of its degrees, and put
$M(\lambda)=\lambda^{2}I-\lambda A+\Delta-I$.

(1) \emph{If $M(\lambda)g=0$ for a real $\lambda>1$ and a vector $g\ne0$, then $\lambda$ is an
eigenvalue of $B$.} A direct computation shows that $f(i\to j)=\lambda g(j)-g(i)$ satisfies
$Bf=\lambda f$. And $f=0$ would give $g(i)=\lambda g(j)=\lambda^{2}g(i)$ along every edge, hence
$g=0$; at an isolated vertex, the equation $M(\lambda)g=0$ is $(\lambda^{2}-1)g=0$.

(2) \emph{If $Bf=\lambda f$ with $f\ne0$ and $\lambda^{2}\ne1$, then $M(\lambda)g=0$ for
$g(j)=\sum_{k\sim j}f(j\to k)$, and $g\ne0$.} The eigenvalue equation at $(i\to j)$ and at
$(j\to i)$ reads $\lambda f(i\to j)=g(j)-f(j\to i)$ and $\lambda f(j\to i)=g(i)-f(i\to j)$.
Eliminating $f(j\to i)$ gives $(\lambda^{2}-1)f(i\to j)=\lambda g(j)-g(i)$. This shows that
$g\ne0$. Summing it over the neighbours $j$ of $i$ gives $M(\lambda)g=0$.

(3) \emph{$B$ has a real eigenvalue $\lambda_1\ge2-d_J/v$.} The matrix
$M(2)=2(\Delta-A)+(3I-\Delta)$ is positive semidefinite, while
$\mathbf{1}^{\!\top}M(\lambda)\mathbf{1}=(\lambda-1)(v\lambda-2v+d_J)$ vanishes at
$\lambda^{*}=2-d_J/v>1$. The smallest eigenvalue of the symmetric matrix $M(\lambda)$ is continuous
in $\lambda$, at most $0$ at $\lambda^{*}$ and at least $0$ at $2$. So $M(\lambda_1)$ is singular for
some $\lambda_1\in[\lambda^{*},2]$, and (1) applies.

(4) \emph{Every non-real eigenvalue $\lambda$ of $B$ satisfies $|\lambda|\le\sqrt2$.} With $g$ as
in~(2), $g^{*}M(\lambda)g=0$ is a quadratic equation $a\lambda^{2}-b\lambda+c'=0$ with real
coefficients $a=g^{*}g>0$, $b=g^{*}Ag$ and $c'=g^{*}(\Delta-I)g\le2a$. Its roots are $\lambda$ and
$\bar\lambda$, so $|\lambda|^{2}=c'/a\le2$.

Now $\tr B^{\ell}$ is the sum of the $\ell$th powers of the $2|E(J)|\le3v$ eigenvalues of $B$. For
even $\ell$ the real eigenvalues contribute nonnegative terms, one of which is at least
$(2-d_J/v)^{\ell}$, and the non-real ones contribute at least $-3v\,2^{\ell/2}$ in total.
\end{proof}

Facts (2) and (4) are part of the Ihara--Bass theory~\cite{Bass,KS}. The test vector $\mathbf{1}$
in~(3) is from~\cite{KNT}. With this test vector the argument needs no assumption of regularity
or expansion.

\begin{lemma}[Cycle weight]\label{lem:weight}
Suppose that $d_J<v$, and let $\ell_0\le L$ be a positive integer. The total weight
$\mu:=\sum_{C\in\mathcal{C}}2^{-|C|}$ satisfies
\[
  \Bigl(1+\tfrac32\,L^{2}\,2^{-s}\Bigr)\,\mu\;\ge\;
  \Bigl(1-\frac{L\,d_J}{2v}-3v\,2^{-\ell_0/2}\Bigr)\sum_{\ell}\frac1{2\ell},
\]
where the sum is over the even $\ell$ with $\ell_0\le\ell\le L$.
\end{lemma}

\begin{proof}
$\tr B^{\ell}$ is the number of closed walks $u_0u_1\cdots u_\ell=u_0$ in $J$ with
$u_{i+1}\ne u_{i-1}$ for all $i$ modulo $\ell$. A cycle of length $\ell$ gives $2\ell$ of them; let
$n_\ell$ be the number of the others. Then
\[
  T:=\sum_{\ell\le L}\frac{\tr B^{\ell}}{2\ell\,2^{\ell}}=\mu+\sum_{\ell\le L}\frac{n_\ell}{2\ell\,2^{\ell}} .
\]
Consider a walk counted by $n_\ell$, and let $j$ be the least index with
$u_j\in\{u_0,\dots,u_{j-1}\}$, say $u_j=u_i$. Then $j<\ell$, the walk $C=u_iu_{i+1}\cdots u_j$ is a
cycle, of length $q$ say, and $u_ju_{j+1}\cdots u_\ell u_1\cdots u_i$ is a non-backtracking closed
walk of length $\ell-q>0$ from $u_i$ to itself. It contains a cycle, so $\ell-q>D>s$. The original
walk is determined by $C$, by the vertex $u_i$ and a direction on $C$ ($2q$ choices), by the
complementary walk (at most $3\cdot2^{\ell-q-s-1}$ choices, by Lemma~\ref{lem:walks}), and by the
position of $u_0$ on the latter (at most $\ell$ choices). Hence
\[
  \sum_{\ell\le L}\frac{n_\ell}{2\ell\,2^{\ell}}
  \;\le\;\sum_{\ell\le L}\ \sum_{C\in\mathcal{C}}\tfrac32\,|C|\,2^{-|C|-s}
  \;\le\;\tfrac32\,L^{2}\,2^{-s}\,\mu ,
\]
and $T\le(1+\frac32L^{2}2^{-s})\mu$. On the other hand, all terms of $T$ are nonnegative. For even
$\ell$ with $\ell_0\le\ell\le L$, Lemma~\ref{lem:spectrum} gives
\[
  2^{-\ell}\tr B^{\ell}\;\ge\;\Bigl(1-\frac{d_J}{2v}\Bigr)^{\ell}-3v\,2^{-\ell/2}
  \;\ge\;1-\frac{L\,d_J}{2v}-3v\,2^{-\ell_0/2}.\qedhere
\]
\end{proof}

\paragraph{The second moment.}
Let $\xi$ be uniform on $\Z(\Gamma)$. For $C\in\mathcal{C}$ let $I_C$ be the event that $C$ is a
component of $\xi$. As the vertices of $C$ have degree $3$ in $\Gamma$, this is the event
$E(C)\subseteq\xi$. It is also the event that $\cO_{V(C)}$ holds and
$\xi\cap E(\Gamma[V(C)])=E(C)$. Put
\[
  a_C=2^{-|C|},\qquad p_C=\PP(I_C),\qquad X_C=\frac{a_C}{p_C}\,\mathbf{1}_{I_C},\qquad
  Z=\sum_{C\in\mathcal{C}}X_C .
\]
Because of the factor $a_C/p_C$, the mean of $X_C$ is the combinatorial weight $a_C$ for every
value of $p_C$. The event $I_C$ is given by $|C|$ linear conditions on $\xi\in\Z(\Gamma)$,
which are satisfied by $\xi=C$. So $p_C\ge a_C$. Hence $0\le X_C\le1$, and $\EE Z=\mu$. If
$Z\ne0$, then $\xi$ contains a cycle of $J$.

Let $C\ne C'$ be in $\mathcal{C}$. If they share a vertex, then $I_C\cap I_{C'}=\emptyset$. If they
are disjoint, Lemma~\ref{lem:link}, applied to the regions $V(C)$ and $V(C')$ of $\Gamma$, gives
$\PP(I_C\cap I_{C'})=2^{\ell(C,C')-1}p_Cp_{C'}$ and therefore
\begin{equation}\label{eq:pair}
  \EE(X_CX_{C'})=a_Ca_{C'}\,2^{\ell(C,C')-1},
\end{equation}
where $\ell(C,C')$ is the number of links between $V(C)$ and $V(C')$ in $\Gamma$.

\begin{lemma}[Second moment]\label{lem:moment}
Suppose that $\ell(C,C')\le b$ for all disjoint $C,C'\in\mathcal{C}$. Let $W_2$ be the sum of
$a_Ca_{C'}$ over the ordered pairs of disjoint cycles of $\mathcal{C}$ with exactly two links. Then
$W_2\le\mu^{2}$ and
\[
  \EE Z^{2}\;\le\;\mu^{2}+W_2+\bigl(1+2^{b}L^{3}2^{-s}\bigr)\mu .
\]
\end{lemma}

\begin{proof}
Call a disjoint pair $C,C'\in\mathcal{C}$ \emph{exceptional} if $\ell(C,C')\ge3$. Fix $C$, and let
$\mathcal{E}$ be a maximal family of pairwise disjoint exceptional partners of $C$. By
Lemma~\ref{lem:exc}, $|\mathcal{E}|\le|\delta V(C)|\le L$. So the union of their vertex sets has
at most $L^{2}$ vertices, and every exceptional partner of $C$ meets this union.
By~\eqref{eq:load}, the sum of $a_{C'}$ over the exceptional partners $C'$ of $C$ is at most
$L^{3}2^{-s}$. With~\eqref{eq:pair},
\[
  \sum_{C'\ \mathrm{exceptional}}\EE(X_CX_{C'})\;\le\;a_C\,2^{b}L^{3}\,2^{-s}.
\]
By~\eqref{eq:pair}, a disjoint pair with at most one link has $\EE(X_CX_{C'})\le a_Ca_{C'}$, and
a pair with two links has $\EE(X_CX_{C'})=2a_Ca_{C'}$. Intersecting pairs contribute nothing, and
$\EE X_C^{2}\le\EE X_C=a_C$. Sum over all ordered pairs.
\end{proof}

\begin{lemma}[Links through trees]\label{lem:treelinks}
Let $8\le D\le L$, and let $t$ be an integer with $1\le t\le D/4$. Let $C,C'\in\mathcal{C}$ be
disjoint, and put $U'=V(C)\cup V(C')$. Consider the links between $V(C)$ and $V(C')$ of two
kinds: the edges; and the components $Y$ of $\Gamma-U'$ that avoid $U$, induce a tree, and have
at most $t$ edges to $U'$. There are at most $289(L/D)^{2}$ such links.
\end{lemma}

\begin{proof}
Let $Y$ be a link of the second kind, with $d_Y\le t$ edges to $U'$. All its vertices have degree
$3$, so $3|Y|=2(|Y|-1)+d_Y$, that is, $|Y|=d_Y-2$. So there is a path from $C$ to $C'$ through $Y$
with at most $t$ edges. For a link that is an edge, take this edge as the path. These paths lie
in $J$. They are pairwise disjoint, because every vertex of $U'$ lies on exactly one edge outside
its own cycle. Divide each of $C$ and $C'$ into arcs of at most $\lfloor D/8\rfloor$ consecutive
vertices. At most $17L/D$ arcs are necessary for each cycle, because $\lfloor D/8\rfloor\ge D/16$
and $L\ge D$. If two of the paths had their ends in the same arc of $C$ and in the same arc of
$C'$, they would form, together with parts of these two arcs, a cycle of $J$ of length at most
$2t+2\lfloor D/8\rfloor<D$. So there are at most $(17L/D)^{2}$ paths.
\end{proof}

\subsection{Proof of Theorem~\ref{thm:forest}}

Let $\Gamma$ and $P$ be as in Lemma~\ref{lem:reduce}, let $\xi$ be uniform on $\Z(\Gamma)$, and let
$r$ be large. All $o(\cdot)$ and $O(\cdot)$ refer to $r\to\infty$ and are uniform in $G$ and $W$.
Put
\begin{gather*}
  x=\log r,\qquad y=\log x,\qquad z=\log y,\\
  D=\lfloor x/10\rfloor,\qquad L=\lfloor x\sqrt y\rfloor,\qquad
  s=\lfloor(D-1)/2\rfloor,\qquad\varepsilon=20/z .
\end{gather*}
Recall that $|P|\le2r/x^{2}$ and $c\le x/100$. Let $\mathcal{F}$ be the event that
$\xi\cap E(\Gamma-P)$ is a forest. We assume that
\begin{equation}\label{eq:assume}
  \PP(\mathcal{F})>\tfrac12+\varepsilon
\end{equation}
and derive a contradiction. By Lemma~\ref{lem:reduce}(ii) this proves Theorem~\ref{thm:forest}.

\paragraph{Removing the short cycles.}
Let $\mathcal{Q}$ be a maximal family of pairwise disjoint cycles of $\Gamma-P$ of length at most
$D$, loops and pairs of parallel edges included, and let $S$ be the union of their vertex sets. For
$Q\in\mathcal{Q}$ the set $V(Q)$ is a cyclic region of $\Gamma$ with $|\delta V(Q)|\le|Q|\le D$,
because its vertices have degree $3$. On $\mathcal{F}$ the set $\xi\cap E(\Gamma[V(Q)])$ is a forest
for every $Q$. So Lemma~\ref{lem:smallcuts} and~\eqref{eq:assume} give
$|\mathcal{Q}|<(D+1)2^{D-1}/\varepsilon\le\sqrt r$. Hence $|S|\le D\sqrt r$.

Put $U=P\cup S$ and $J=\Gamma-U$. By the maximality of $\mathcal{Q}$, $J$ is a simple graph of
girth greater than $D$. So we are in the setting of Section~\ref{sec:moderate}. The number $v$ of
vertices of $J$ satisfies $v=(2+o(1))r$, and
\begin{equation}\label{eq:deficit}
  d_J\;\le\;3\bigl(|P|+|S|\bigr)\;\le\;\frac{7r}{x^{2}} .
\end{equation}

\paragraph{The weight of the cycles.}
We claim that
\begin{equation}\label{eq:mu}
  \mu\;\ge\;\bigl(\tfrac18-o(1)\bigr)\ln y .
\end{equation}
This is Lemma~\ref{lem:weight} with $\ell_0=\lceil6x\rceil$. Indeed $2^{-s}\le8\,r^{-1/20}$ and
$L^{2}\le x^{2}y$, so $L^{2}2^{-s}=o(1)$. By~\eqref{eq:deficit}, $L\,d_J/(2v)=O(\sqrt y/x)$. Also
$3v\,2^{-\ell_0/2}\le3v\,r^{-3}=o(1)$. The sum of $1/(2\ell)$ over the even $\ell$ in $[6x,L]$ is
$(\frac14-o(1))\ln(L/6x)=(\frac18-o(1))\ln y$.

On $\mathcal{F}$ no cycle of $J\subseteq\Gamma-P$ is contained in $\xi$, so $Z=0$. By~\eqref{eq:assume},
\begin{equation}\label{eq:Zzero}
  \PP(Z=0)>\tfrac12+\varepsilon .
\end{equation}

\begin{lemma}[Few links]\label{lem:links}
$\ell(C,C')\le c+o(x)$ for all disjoint $C,C'\in\mathcal{C}$.
\end{lemma}

\begin{proof}
Let $U'=V(C)\cup V(C')$. Every vertex of $U'$ lies on exactly one edge outside its own cycle, so at
most $2L$ edges leave $U'$. At most $c$ components of $\Gamma-U'$ meet $P$, because
$P\cap U'=\emptyset$ and so each component of $\Gamma[P]$ lies in one of them. Let $Y$ run over the
other components and let $d_Y$ be the number of edges between $Y$ and $U'$; then
$\sum_Yd_Y\le2L$. Put $t_0=\lfloor y/2\rfloor$.

\step{Large cuts.} At most $2L/t_0$ components have $d_Y>t_0$.

\step{Small cuts, with a cycle.} The components $Y$ with $d_Y\le t_0$ that contain a cycle are
pairwise disjoint cyclic regions inside $\Gamma-P$ with $|\delta Y|\le t_0$. On $\mathcal{F}$ each
$\xi\cap E(\Gamma[Y])$ is a forest, so by Lemma~\ref{lem:smallcuts} and~\eqref{eq:assume} there are
fewer than $(t_0+1)2^{t_0-1}/\varepsilon$ of them.

\step{Small cuts, trees.} Let $Y$ induce a tree, with $d_Y\le t_0$. No $Q\in\mathcal{Q}$ meets $Y$,
since $Q$ is a connected subgraph of $\Gamma-U'$ and would lie in $Y$. Hence $Y$ avoids $U$. By
Lemma~\ref{lem:treelinks} with $t=t_0$, at most $289(L/D)^{2}$ links are such components or edges.

Altogether
\[
  \ell(C,C')\;\le\;c+\frac{2L}{t_0}+\frac{(t_0+1)2^{t_0-1}}{\varepsilon}+289\Bigl(\frac{L}{D}\Bigr)^{2}
  =c+O\Bigl(\frac{x}{\sqrt y}\Bigr)+O\bigl(yz\sqrt x\bigr)+O(y)=c+o(x).\qedhere
\]
\end{proof}

\paragraph{The contradiction.}
By Lemma~\ref{lem:links} we may use Lemma~\ref{lem:moment} with $b=c+o(x)$. As $c\le x/100$,
\[
  2^{b}L^{3}2^{-s}=2^{\,c-x/20+o(x)}=o(1).
\]
So $\EE Z^{2}\le2\mu^{2}+(1+o(1))\mu$. By Lemma~\ref{lem:second} and~\eqref{eq:mu},
\[
  \PP(Z=0)\;\le\;\frac12+\frac{1+o(1)}{4\mu}\;\le\;\frac12+\frac{2+o(1)}{\ln y}\;<\;\frac12+\frac{20}{z}
\]
for large $r$, since $\ln y=z\ln2$. This contradicts~\eqref{eq:Zzero}, and
Theorem~\ref{thm:forest} is proved.

\subsection{The forest input and the leading term}\label{sec:leading}

Later sections use the theorem for one cycle in a multigraph, with a fixed error.

\begin{proposition}[Forest input]\label{prop:input}
For every $\delta\in(0,\half)$ there is an integer $R_\delta\ge2^{100}$ with the following property.
Let $H$ be a connected multigraph of rank $\rho\ge R_\delta$, and let $W$ be a cycle of $H$ with at
least three and at most $\rho/\log^{2}\rho$ edges. Then $p(H,W)\le\half+\delta$.
\end{proposition}

\begin{proof}
Replace every edge of $H$ outside $W$ by a path with three edges. The new graph $H'$ is simple and
connected, it contains $W$, and it has rank $\rho$. The even subgraphs of $H'$ correspond to those
of $H$, and one has a forest outside $W$ if and only if the other has. Apply
Theorem~\ref{thm:forest} to $H'$ and $W$. Choose $R_\delta$ so large that $\rho\ge R_\delta$ is in
the range of the theorem, $\log\rho\ge100$, and $20/\lgk3\rho\le\delta$.
\end{proof}

\begin{lemma}[Cubic multigraphs]\label{lem:cubic}
For every $r\ge2$ there is a connected cubic multigraph of rank $r$ with exactly $\psi(r)$ cycles.
\end{lemma}

\begin{proof}
The numbers of cycles of the graphs of rank $r$ are integers below $2^{r}$, so some graph $G_0$
attains $\psi(r)$. If $G_0$ is not connected, identify a vertex of one component with a vertex of
another. This changes neither the rank nor the set of cycles. Now delete a vertex of degree one, or
replace a vertex of degree two and its two edges by one edge, as long as this is possible. These
steps keep the rank and the number of cycles. The result is a connected multigraph with all
degrees at least three, because $r\ge2$.

Let $v$ be a vertex of degree $d\ge4$. Replace $v$ by a path $v^{1}v^{2}\cdots v^{d-2}$, and
distribute the $d$ edge ends at $v$ along this path: two at $v^{1}$, one at each of
$v^{2},\dots,v^{d-3}$, and two at $v^{d-2}$. The new multigraph has $d-3$ more vertices and $d-3$
more edges, so its rank is $r$. A cycle through $v$ uses two edge ends at $v$. Together with the
part of the path between these two ends it is a cycle of the new multigraph. Different cycles give
different cycles, and the cycles that avoid $v$ remain. So the number of cycles does not decrease,
and it cannot exceed $\psi(r)$. Repeat this for every vertex of degree at least four.
\end{proof}

\begin{corollary}[Leading term]\label{cor:cycles}
There is $r_1$ such that every graph of rank $r\ge r_1$ has at most
$\bigl(\half+21/\lgk3r\bigr)2^{r}$ cycles. Consequently $\psi(r)\sim2^{r-1}$.
\end{corollary}

\begin{proof}
Let $r$ be large, and let $G$ be a connected cubic multigraph of rank $r$ with $\psi(r)$ cycles
(Lemma~\ref{lem:cubic}). Let $W$ be a shortest cycle of $G$, of length $g$. A breadth-first search
from a vertex meets no vertex twice up to distance $\lfloor(g-1)/2\rfloor$, so $G$ has at least
$2^{(g-2)/2}$ vertices. As $G$ has $2r-2$ vertices,
\begin{equation}\label{eq:girthbound}
  g\;\le\;2\log(2r)+2 .
\end{equation}
If $g<3$, subdivide an edge of $W$ twice; this changes neither $c(G)$ nor $p(G,W)$. Now $W$ has
between $3$ and $r/\log^{2}r$ edges. Subdivide the edges outside $W$ as in the proof of
Proposition~\ref{prop:input}. Then Theorem~\ref{thm:forest} gives
$p(G,W)\le\half+\delta_r$ with $\delta_r=20/\lgk3r$. Let $e$ be an edge of $W$. By
Lemma~\ref{lem:edge},
\[
  \psi(r)=c(G)\;\le\;c(G-e)+\bigl(\tfrac12+\delta_r\bigr)2^{r-1}
  \;\le\;\psi(r-1)+\bigl(\tfrac12+\delta_r\bigr)2^{r-1}.
\]
Put $y_r=\psi(r)2^{-r}-\half$. Then $y_r\le\half y_{r-1}+\half\delta_r$ for all large $r$. The
numbers $\delta_r$ decrease, and $y_r\le\half$. So $y_r\le2^{-j}y_{r-j}+\delta_{r-j}$ for
$j\le r/2$. With $j=\lfloor r/2\rfloor$ this is at most $21/\lgk3r$ when $r$ is large. Finally
$\psi(r)>2^{r-1}$ for $r\ge3$~\cite{ES}; the rings of Section~\ref{sec:rings} also show this.
\end{proof}
